\section{Introduction}\label{sec:intro}

Many decisions in engineering and in statistics are made by simulating a finite set of alternatives and comparing them on more than one criterion. An inventory manager simulates reorder policies and weighs their operating cost against how often customers wait; a methodologist simulates confidence intervals and weighs their coverage against their length; a health agency simulates vaccination strategies and weighs several immunological responses \citep{kone2023}. Each run returns a noisy vector of performance measures, and the object of interest is the Pareto set, the alternatives not dominated in expectation. The usual practice is to fix a number of replications and read the Pareto set off the averages, which can make a good alternative look dominated through sampling error alone and does not show which conclusions the replications support.

Two features of the practice make formal inference difficult. First, analysts look at preliminary results, add replications where a comparison seems close and stop when the picture looks stable; classical intervals computed after such data-dependent stopping lose their coverage \citep{johari2022,ramdas2023}. Second, the conclusions concern many alternatives and objectives at once. Confidence sequences address the first difficulty: they contain a mean simultaneously at all sample sizes with a prescribed probability, so that stopping at any data-dependent time is permitted \citep{darling1967,howard2021}. For bounded observations, the betting confidence sequences of \citet{waudbysmith2024} are nonparametric and adapt to the variance. The second difficulty can be met by running one confidence sequence for every alternative and objective and splitting the error level among them. The product of the intervals of an alternative is then a confidence box for its mean vector, and statements about dominance can be read from the boxes.

This paper grew out of a failure. In a first version of our software, an alternative was discarded when the box of another alternative showed that the latter was better up to a tolerance~$\varepsilon$ in every objective. A positive tolerance is needed: without it two alternatives with equal means can never be separated and the procedure does not stop. Every certificate issued by that rule was correct, in the sense that each discarded alternative was within the tolerance of the alternative that discarded it. We nevertheless found configurations in which the rule discarded every Pareto-optimal alternative and reported, with full confidence, an empty set. The reason is that dominance up to a tolerance is not transitive: a chain, or a cycle, of individually correct statements does not compose into a correct conclusion. The same phenomenon is known in evolutionary multiobjective optimization, where archives based on $\varepsilon$-dominance were redesigned to avoid deterioration \citep{laumanns2002}. We had not expected it to arise in an inferential procedure in which every individual statement is protected by a simultaneous guarantee.

The remedy we propose separates the two roles that the tolerance was playing. Alternatives are discarded only on evidence of exact Pareto dominance, a relation that is transitive and acyclic on confidence boxes; the tolerance is used only to decide when a surviving alternative may be accepted, namely when no other alternative can improve on it by more than $\varepsilon$ in every objective. This is the accuracy criterion of \citet{auer2016} for Pareto front identification. The contribution of the paper is to obtain it with certificates that are valid uniformly over time, for bounded outcomes with arbitrary dependence among the objectives of one replication, and to make the consequences of the design choices explicit. Specifically:
\begin{enumerate}[leftmargin=*,itemsep=0pt]
\item we prove that discarding by certified $\varepsilon$-dominance can remove the whole Pareto set, with three objectives through a cycle and with two objectives through the tie-breaking that such a rule requires (Section~\ref{sec:failure});
\item we give a procedure whose certificates, read at any data-dependent time, contain the Pareto set, accept only $\varepsilon$-Pareto-optimal alternatives and provide a certified $\varepsilon$-cover of the Pareto set, with probability at least $1-\alpha$ (Theorem~\ref{thm:validity});
\item we show that the natural allocation rule, which evaluates only alternatives whose status is undecided, can stall, and that evaluating both members of every open comparison yields a deterministic bound on the number of evaluations (Theorem~\ref{thm:termination});
\item we evaluate the procedure in a simulation study and in two applications where the truth is known exactly, and we provide an R implementation \citep{paretoinfer}.
\end{enumerate}

We examine three hypotheses. H1: discarding by $\varepsilon$-dominance loses Pareto-optimal alternatives in configurations with near-ties, whereas discarding by exact dominance does not. H2: the certificates of the proposed procedure keep the probability of any false statement below $\alpha$ under optional stopping and adaptive allocation, whereas fixed-sample intervals recomputed after every batch, and confidence sequences without a multiplicity correction, do not. H3: allocating evaluations to open comparisons reduces the number of evaluations relative to uniform allocation, increasingly so as the number of alternatives grows.

\paragraph{Related work.} Multiobjective ranking and selection, reviewed by \citet{hunter2019}, has concentrated on allocating a fixed simulation budget so that the probability of a correct Pareto set is high, typically under normality; examples are the MOCBA procedure of \citet{lee2010} and the SCORE allocations of \citet{feldman2018}. These procedures are designed for a budget fixed in advance and do not offer guarantees after data-dependent stopping. In the bandit literature, Pareto set identification with fixed confidence was studied by \citet{auer2016}, who introduced the $\varepsilon$-accuracy criterion used here and derived problem-dependent lower and upper bounds, and later by \citet{kone2023} and \citet{crepon2024}, who obtained algorithms with near-optimal or optimal sample complexity under sub-Gaussian or Gaussian assumptions; \citet{ararat2023} treated general ordering cones. These works guarantee the output at the time fixed by their own stopping rule, under independent objectives or known distributional families; the $\varepsilon$-PAL algorithm of \citet{zuluaga2016} uses a Gaussian process model over a large design space. Our procedure is less efficient than the best of these methods when their assumptions hold, and it does not exploit common random numbers. In return its statements are nonparametric, remain valid whenever the analyst stops, including after a deadline or a budget overrun, and leave an alternative undecided when the evidence is insufficient.

Section~\ref{sec:problem} states the problem, Section~\ref{sec:failure} the failure, Section~\ref{sec:method} the procedure and its guarantees, Sections~\ref{sec:simulation} and~\ref{sec:applications} the simulation study and the applications, and Section~\ref{sec:discussion} the limitations. Proofs, additional results and the code are in the Supplementary Material.
